% ============================================================
%  PREAMBOLO CONDIVISO PER LE VERIFICHE
%  (tutto tranne il font, caricato a monte da chi include questo
%  file: preambolo-verifica.tex oppure preambolo-verifica-dsa.tex)
%
%  Versione ridotta di ../latex/preamble.tex: qui mancano di
%  proposito tutte le cose pensate per gli appunti (mdanimation,
%  riassunto, provatu, QR nel margine, modalitaesercizi/rigaesercizi,
%  gli ambienti exercise/expsolution con link "vai alla soluzione")
%  perché in una verifica non servono mai.
% ============================================================

% ============================================================
%  LINGUE
% ============================================================
\usepackage{polyglossia}
\setmainlanguage{italian}

% ============================================================
%  MATEMATICA
% ============================================================
\usepackage{amsmath}
\usepackage{amssymb}
\usepackage{amsfonts}
\usepackage{amsthm}
\usepackage{mathtools}
\usepackage{cancel}

% Spaziatura display math
\setlength{\abovedisplayskip}{4pt}
\setlength{\belowdisplayskip}{4pt}
\setlength{\abovedisplayshortskip}{2pt}
\setlength{\belowdisplayshortskip}{2pt}

% Comandi matematici
\newcommand{\R}{\mathbb{R}}
\newcommand{\tonde}[1]{\left(#1\right)}
\newcommand{\quadre}[1]{\left[#1\right]}
\newcommand{\graffe}[1]{\left\lbrace #1 \right\rbrace}

% ============================================================
%  UNITÀ DI MISURA
% ============================================================
\usepackage{siunitx}
\sisetup{
  detect-all,
  output-decimal-marker = {,}
}

% ============================================================
%  LAYOUT E GEOMETRIA
% ============================================================
\usepackage[centering, includeheadfoot]{geometry}
\geometry{
  paperwidth  = 170mm,
  paperheight = 240mm,
  top         = 1.5cm,
  bottom      = 1.5cm,
  inner       = 1.5cm,
  outer       = 1.5cm,
  twoside,
}

% Spaziatura testo
\usepackage{setspace}
\setstretch{1.18}
\usepackage{parskip}
\setlength{\parindent}{0pt}
\setlength{\parskip}{0.25\baselineskip plus 1pt minus 1pt}

\raggedbottom

% ============================================================
%  INTESTAZIONI E PIÈ DI PAGINA
% ============================================================
\usepackage{fancyhdr}
\pagestyle{fancy}
\fancyhf{}
\fancyfoot[C]{\thepage}
\renewcommand{\headrulewidth}{0pt}
\renewcommand{\footrulewidth}{0pt}

% ============================================================
%  TITOLI
% ============================================================
\usepackage{titlesec}

\titleformat{\chapter}[block]
  {\normalfont\LARGE\bfseries\color{cyan}}
  {\rule{\textwidth}{.5pt}\\[0.1em]\centering}
  {1em}
  {\huge\scshape\centering}
  [\vspace{-0.2em}\rule{\textwidth}{0.5pt}]

\titleformat{\section}[block]
  {\normalfont\Large\bfseries\color{cyan}}{\thesection}{1em}{}

\titlespacing*{\chapter} {0pt}{-30pt}{8pt}
\titlespacing*{\section} {0pt}{0.8\baselineskip}{0.3\baselineskip}

% ============================================================
%  LISTE
% ============================================================
\usepackage[inline]{enumitem}
\setlist{itemsep=0pt, parsep=0pt, topsep=3pt, partopsep=0pt}
% Itemsep più generoso per gli esercizi delle verifiche (aumentato
% rispetto al default: qui lo spazio serve per far svolgere i calcoli).
\setlist[enumerate]{itemsep=14pt}
\newlist{sceltamultipla}{enumerate}{1}
\setlist[sceltamultipla,1]{%
  label=\textcolor{gray}{\fbox{\textbf{\footnotesize\makebox[1.2em]{\Alph*}}}},
  leftmargin=*,
  itemsep=1pt,
  topsep=2pt
}

% ============================================================
%  COLORI
% ============================================================
\usepackage{xcolor}
\definecolor{ciano}{RGB}{0, 180, 200}
\definecolor{gold}  {RGB}{255, 215, 0}

% ============================================================
%  BOX: MDFRAMED
% ============================================================
\usepackage{mdframed}

\newenvironment{mdexam}[1][]
  {\begin{mdframed}[topline=false,bottomline=false,rightline=false,
    linecolor=magenta,linewidth=2,frametitlefont=\itshape]
   \textbf{Allenamento per la verifica. }#1\ignorespaces}
  {\end{mdframed}}

% Regola d'oro della verifica: da richiamare in cima a ogni verifica.
\newcommand{\regoladoro}{%
  \begin{mdframed}[linewidth=2]
    \textbf{Regola d'oro della verifica.} \textit{Una verifica in
    bianco vale 4, una verifica perfetta ma copiata vale 2.}
  \end{mdframed}%
}

% ============================================================
%  TESTO E TIPOGRAFIA
% ============================================================
\usepackage[normalem]{ulem}
\usepackage{soul}
\setulcolor{gray}
\usepackage{paracol}

% ============================================================
%  PACCHETTI GRAFICI E VARI
% ============================================================
\usepackage{graphicx}
\usepackage{array}
\usepackage{booktabs}
